\bookchapter{الفصل 8: الاتصالات: الرسائل النصية والمكالمات والتواصل الاجتماعي}

\chapterintro{الاتصالات هي المهمة الأصلية والجوهرية للهاتف المحمول. لقد تطورت الآن من المكالمات والرسائل النصية البسيطة إلى نظام بيئي معقد يشمل المراسلة الفورية والشبكات الاجتماعية والبريد الإلكتروني.}

\booksection{أولاً: الميزات المتقدمة لتطبيق الهاتف}

تجاوز تطبيق الهاتف في الهواتف الحديثة لوحة الاتصال البسيطة، فهو يضم أدوات قوية مثل التسجيل وحظر المكالمات المزعجة والتعرف الذكي.

\booksubsection{تسجيل المكالمات (مع مراعاة الامتثال القانوني):}

\begin{itemize}
  \item الفائدة: توثيق محتوى المكالمات المهمة لمحاضر العمل أو المقابلات أو التزامات العملاء أو حفظ الأدلة عند النزاعات.
  \item تنبيه قانوني: داخل الصين، يُشترط لاستخدام تسجيل المكالمات كدليل الحصول على موافقة الطرف الآخر، أو استخدامه للدفاع عن الحقوق والمصالح المشروعة.
  \item التسجيل أحادي الطرف مقابل موافقة الطرفين: قانونيًا، التسجيل الذي يُحصل عليه «بدون موافقة الطرف الآخر» ليس باطلاً بحد ذاته. فحسب التفسيرات القضائية ذات الصلة للمحكمة الشعبية العليا، يُمكن تقديم التسجيل كدليل طالما أن طريقة الحصول عليه لم تنتهك بجسامة حقوق الآخرين المشروعة، ولم تخالف أحكامًا قانونية صارمة، ولم تتعارض مع الآداب العامة. لكن بعض محاكم المحافظات تكون أكثر صرامة في قبول «التسجيل السري» في القضايا المدنية؛ أما القضايا الجنائية فعادةً تخضع لمراجعة قانونية أكثر صرامة.
  \item الاختلافات الإقليمية: محاكم بكين وشنغهاي وغوانغدونغ أكثر تساهلاً مع التسجيل أحادي الطرف كدليل مدني؛ أما التسجيلات التي تتعلق بأسرار تجارية أو خصوصية شخصية فتحتاج إلى حذر خاص في أي محافظة، وقد تُستبعد لانتهاكها حق الخصوصية.
  \item كيفية التنفيذ: يتوفر عادةً في Android زر «تسجيل» واضح في شاشة المكالمة؛ أما نظام iOS فلا يوفر ميزة تسجيل المكالمات الأصلية، لذا تحتاج إلى تطبيق طرف ثالث أو سماعات بلوتوث بها ميزة تسجيل.
\end{itemize}

\begin{tcbwarn}{خطوط حمراء لاستخدام التسجيل}
\begin{itemize}
  \item لا يجوز استخدام التسجيل للابتزاز أو النشر أو الإهانة العلنية للآخرين.
  \item المحتوى الذي يمس خصوصية الغير (مثل تسجيل حديث أشخاص آخرين عن طريق الصدفة) يجب عدم نشره دون ضوابط.
  \item في المواقف المهمة، يُنصح بالإعلان شفهيًا في بداية المكالمة بأن «هذه المكالمة ستُسجَّل»، وهو ما يضمن الامتثال ويزيد من القوة الإثباتية.
\end{itemize}

\end{tcbwarn}

\booksubsection{محتوى متعمق حول التعامل مع المكالمات المزعجة:}

\begin{itemize}
  \item المبدأ: يستعلم تطبيق الهاتف عبر الإنترنت عن قاعدة بيانات ضخمة من الأرقام، ويعرض عند الاتصال «مكالمة مزعجة» أو «مكالمة احتيال» أو اسم الشركة.
  \item كيفية التفعيل: في Android افتح تطبيق الهاتف > الإعدادات > حظر المزعجات؛ في iOS افتح الإعدادات > الهاتف > منع الاتصالات والتعريف. لمسارات التفعيل التفصيلية لكل علامة تجارية والخيارات الموصى بها، راجع جدول المقارنة بين العلامات التجارية في الفصل 13 «أدوات النظام المفيدة».
  \item منصات الوسم الرئيسية: يُعتمد في الصين بشكل رئيسي على ثلاث شركات لبيانات وسم الأرقام — Tencent Mobile Manager و360 Mobile Guard وPhone Helper (Tapir). كما تستدعي شركات الاتصالات (China Mobile / China Unicom / China Telecom) هذه البيانات غالبًا لعرض «اسم صاحب الرقم» ووظيفة «حظر المزعجات». وتتعاون مصانع الهواتف (Huawei وXiaomi وOPPO وvivo) في قواعد الوسم مع مزودي الخدمة المذكورين.
  \item ماذا تفعل عند الوسم الخاطئ: إذا تم وسم رقم عملك العادي بشكل خاطئ كـ«توصيل طعام» أو «توصيل طرود» أو «مكالمة مزعجة»، يمكنك الاعتراض عبر القنوات التالية: Phone Helper: ازور منصة اعتراض الأرقام (haoma.tk.cn) وقدّم إثبات الهوية.
  \item 360: قدّم الطلب عبر وظيفة «اعتراض الرقم» داخل تطبيق 360 Mobile Guard.
  \item Tencent: قدّم الطلب عبر صفحة «اعتراض الرقم» في تطبيق Tencent Mobile Manager المصغّر.
  \item شركات الاتصالات الثلاث الكبرى: يمكن لمشتركي China Mobile إرسال رسالة نصية «KTFSR» إلى 10086 لتفعيل ميزة منع الإزعاج؛ يتابع مشتركو China Telecom حساب «Tianyi Anti-Harassment»؛ ويضبط مشتركو China Unicom الإعدادات في «China Unicom Mobile Manager».
\end{itemize}

المنع المتقدم: فعّل «حظر المكالمات عالية التردد» لدى شركة الاتصالات (مثل قيود China Mobile للمكالمات عالية التردد) لحظر مكالمات الذكاء الاصطناعي الآلية تلقائيًا.

القدرات الشائعة للحظر المدمج في الهاتف: التعرف على المتصل / الحظر الذكي: وسم وحظر المكالمات المزعجة والاحتيالية والإعلانية تلقائيًا حسب قاعدة بيانات الأرقام.

القائمة السوداء / القائمة البيضاء: احظر يدويًا الأرقام المزعجة المعروفة؛ وأضف جهات الاتصال وأرقام التوصيل إلى القائمة البيضاء لتجنب الحظر الخاطئ.

تصفية الرسائل النصية بالكلمات المفتاحية: تُوضع الرسائل التي تحتوي على كلمات مثل «فزت بجائزة» أو «قرض» أو «انقر على الرابط» تلقائيًا في صندوق الحظر.

كتم المكالمات غير المعروفة في iOS: تُكتم مكالمات الأرقام غير الموجودة في جهات الاتصال وتذهب إلى البريد الصوتي دون مقاطعتك؛ مناسبة لمن يتلقى مكالمات مزعجة كثيرة، لكن انتبه لتفويت مكالمات مهمة من أرقام غريبة.

الترابط مع وضع عدم الإزعاج: سمح فقط بمكالمات جهات الاتصال أثناء الليل أو الاجتماعات، مع المؤقت المجدول، وستتحسن جودة النوم فورًا.

كيفية تنسيق التطبيقات العملية: استخدم أولًا الحظر المدمج في النظام، ثم أضف تطبيق National Anti-Fraud Center وخدمات منع الإزعاج لدى شركة الاتصالات؛ ويمكن اختيار أدوات خفيفة مثل Qingzhuo للتنظيف، لكن يكفي تثبيت مجموعة واحدة من تطبيقات الأمان/التنظيف، والتفاصيل في الفصل 13.

\booksubsection{ما هو VoLTE/VoNR وكيفية تفعيله:}

\begin{itemize}
  \item VoLTE (Voice over LTE): تقنية مكالمات صوتية عبر شبكة 4G، لا يعود الاتصال يهبط إلى 2G/3G أثناء المكالمة، ويمكنك التحدث والتصفح في نفس الوقت، وجودة الصوت أوضح من المكالمات التقليدية والاتصال أسرع.
  \item VoNR (Voice over NR): مكالمات صوتية عبر 5G، وهي الجيل التالي من VoLTE. زمن الاستجابة أقل وجودة صوت أفضل، وتدعم مكالمات الفيديو عالية الدقة. يتطلب هاتف 5G + باقة 5G + تغطية 5G محلية في نفس الوقت.
  \item طريقة التفعيل: في iOS فعّل «VoLTE» من الإعدادات > الشبكة الخلوية > خيارات البيانات الخلوية > الصوت والبيانات؛ في Android ابحث عن مفتاح «مكالمة VoLTE عالية الدقة» في الإعدادات > الشبكة المتنقلة > البيانات المتنقلة > متقدمة. عادةً تكون مفعّلة افتراضيًا من شركة الاتصالات؛ وإذا انقطعت الشبكة أثناء المكالمة أو تعذر التصفح، فغالبًا يكون السبب عدم تفعيل VoLTE.
\end{itemize}

\begin{tcbinfo}{معلومة صغيرة}
إذا أصبحت سرعة الإنترنت بطيئة فجأة أثناء المكالمة أو انقطعت، فغالبًا يكون VoLTE/VoNR غير مفعّل وتم الرجوع إلى شبكة 2G/3G. بعد التفعيل يمكنك «التحدث وتصفح الفيديو في نفس الوقت».

\end{tcbinfo}

\booksection{ثانيًا: الرسائل النصية وiMessage/RCS}

لم تُلغَ الرسائل النصية التقليدية، فهي لا تزال قناة موثوقة لاستقبال رموز التحقق وإشعارات البنوك والرسائل الرسمية.

\begin{itemize}
  \item iMessage: رسائل تُرسل عبر الإنترنت بين أجهزة Apple، وتظهر في فقاعات زرقاء. تدعم صورًا ومقاطع فيديو عالية الجودة وإيصالات القراءة، وهي مجانية بالكامل ومشفرة من طرف إلى طرف.
  \item RCS: بروتوكول اتصالات من الجيل التالي تدعمه شركات الاتصالات العالمية وGoogle. وظائفه مشابهة لـ iMessage أو WeChat، لكنها مبنية على رقم الهاتف.
  \item التوافق بين iMessage وRCS: منذ iOS 18، بدأت Apple بدعم معيار RCS، مما يعني أنه يمكن إرسال صور عالية الدقة وإيصالات القراءة وحالة الكتابة بين iPhone وهواتف Android عبر RCS، والتخلص من انقطاع تجربة «الفقاعات الخضراء». لكن بشرط أن يدعم كلا الطرفين RCS من شركتي الاتصالات، وأن يكون إصدار نظام الهاتف حديثًا بما فيه الكفاية.
\end{itemize}

\booksubsection{حالة نشر RCS في الصين:}

\begin{itemize}
  \item China Mobile: نشرت RCS مبكرًا، وتدعم بعض المحافظات خدمة «رسائل 5G» لإرسال رسائل وسائط غنية وبطاقات تفاعلية. يحتاج إلى دعم الجهاز وتفعيل الخدمة في مكتب الشركة.
  \item China Telecom: تقدم الخدمة عبر «رسائل 5G من China Telecom»، وبعض الباقات تدعمها افتراضيًا، أما المستخدمون القدماء فيحتاجون إلى إرسال أمر رسالة نصية محدد للتفعيل.
  \item China Unicom: نشرت رسائل 5G في المدن الرئيسية، لكن تغطية الأجهزة غير متساوية، ولا يزال العديد من طرازات Android يحتاج إلى تفعيل «الرسائل المحسنة» أو «رسائل 5G» يدويًا في تطبيق الرسائل.
  \item التجربة الفعلية: كان انتشار RCS داخل الصين أبطأ من المتوقع، ويرجع ذلك أساسًا إلى هيمنة منظومة WeChat وعدم توحيد دعم الأجهزة الطرفية. لا تزال الاتصالات اليومية تعتمد أساسًا على WeChat، بينما يضطلع RCS أكثر بدور «الإشعارات الرسمية» (مثل إشعارات البنوك والرسائل الحكومية).
\end{itemize}

\begin{tcbwarn}{خطوط حمراء لأمان رموز التحقق}
\begin{itemize}
  \item لا تُرسل رمز التحقق: أي شخص يطلب منك «قراءة رمز التحقق» أو «تحويل الرسالة النصية» قد يكون محتالًا، بمن فيهم من يدّعي أنه «خدمة العملاء» أو «سائق التوصيل» أو «الشرطة/النيابة/المحكمة». رمز التحقق = محفظتك.
  \item الرسائل النصية التصيدية: احذر الروابط في رسائل مثل «انتهت صلاحية ETC» أو «تجميد بطاقة التأمين الصحي» أو «التعويض عن فقدان الطرد»، فالمؤسسات الرسمية لا تطلب رقم البطاقة البنكية أو كلمة المرور عبر روابط الرسائل النصية.
  \item التعرف على الأرقام: عادةً تأتي الرسائل البنكية الرسمية من أرقام قصيرة تبدأ بـ 95 أو من 1069 مع توقيع الشركة؛ وأي «إشعار بنكي» من رقم جوال عادي مكون من 11 رقمًا هو في الغالب احتيال.
  \item النسخ الاحتياطي للتحقق الثنائي: لا تغيّر رقم الهاتف المرتبط بالحسابات المهمة بشكل عشوائي؛ وإذا اضطررت لتغيير الرقم، ففك الارتباط أولًا بجميع الحسابات الحرجة، وإلا فقد يعيد صاحب الرقم الجديد تعيين كلمة المرور الخاصة بك عبر رمز التحقق.
\end{itemize}

\end{tcbwarn}

\booksection{ثالثًا: WeChat/QQ: إدارة تطبيقات التواصل الاجتماعي الشاملة}

تجاوز WeChat وQQ أدوات الاتصالات، ليصبحا «تطبيقًا شاملاً» يجمع التواصل والدفع والحصول على المعلومات.

\booksubsection{النسخ الاحتياطي ونقل سجلات المحادثات:}

\begin{enumerate}
  \item سجّل الدخول إلى WeChat لنسخة الكمبيوتر أو Mac على جهاز الكمبيوتر.
  \item انقر على القائمة في الزاوية اليسرى السفلية > [النقل والنسخ الاحتياطي] > [النسخ الاحتياطي والاستعادة].
  \item اختر «النسخ الاحتياطي لسجلات المحادثة إلى الكمبيوتر»، ثم حدد على الهاتف المحادثات المراد نسخها احتياطيًا.
  \item النقل إلى هاتف جديد: على الهاتف القديم في WeChat [أنا] > [الإعدادات] > [عام] > [النقل والنسخ الاحتياطي لسجلات المحادثات] > [نقل] > [النقل إلى هاتف/جهاز لوحي]، وامسح الرمز للانتهاء.
\end{enumerate}

\booksubsection{طرق متقدمة لتحرير مساحة التخزين:}

\begin{itemize}
  \item المدخل المعتاد: WeChat [أنا] > [الإعدادات] > [عام] > [مساحة التخزين]، يمكنك تنظيف الذاكرة المؤقتة وتنظيف سجلات المحادثات بدقة حسب الوقت ونوع الملف.
  \item التنظيف داخل المحادثة: ادخل محادثة معينة > «...» في الزاوية اليمنى العلوية > [سجل المحادثات] > [إدارة مساحة التخزين]، يمكنك تنظيف الصور ومقاطع الفيديو والملفات في محادثة محددة مع الاحتفاظ بالنص.
  \item تنظيف الصور/الفيديو الأصلية: تشغل الصور والمقاطع الأصلية عالية الدقة الواردة في المحادثات الجماعية أكبر مساحة، ويمكن حذفها دفعة واحدة بتحديد «الأصلية» في [إدارة مساحة التخزين].
  \item إيقاف التنزيل التلقائي: [الإعدادات] > [عام] > [الصور والفيديو والملفات والمكالمات]، أوقف «التنزيل التلقائي» لتجنب تخزين ملفات المجموعات الكثيرة في الخفاء عبر WiFi.
  \item ذاكرة التطبيقات المصغرة: في [اكتشف] > [التطبيقات المصغرة] اضغط مطولًا على التطبيقات غير المستخدمة لحذفها وتحرير الذاكرة المحلية.
  \item النسخ الاحتياطي المضغوط لسجلات المحادثات: عند النسخ الاحتياطي من WeChat على الكمبيوتر حدد «ضغط» لتقليل سجلات المحادثات من عدة جيجابايت إلى 30\%-50\% من الحجم الأصلي.
\end{itemize}

\booksubsection{حيل إدارة المجموعات:}

\begin{itemize}
  \item إعلان المجموعة: ينشره المالك/المشرف من «...» في الزاوية اليمنى العلوية > [إعلان المجموعة]، مع تنبيه قوي لجميع الأعضاء، مناسب للإشعارات المهمة. يظهر تلقائيًا عند انضمام عضو جديد.
  \item مهام المجموعة: يمكن للمالك تعيين الرسالة المثبتة كـ«مهمة مجموعة»، فتظهر إشارة حمراء فوق رؤوس جميع الأعضاء وتختفي تلقائيًا بعد الإكمال. مناسبة لتسجيل الفعاليات وجمع المهام.
  \item سلسلة التسجيل الجماعي: اكتب «\#سلسلة» في مربع الإدخال لبدء سلسلة، وهي مناسبة لطلب الوجبات وحصر التسجيلات. تُجمَّع نتائج السلسلة تلقائيًا في أعلى محادثة المجموعة، لتجنب إغراق الشاشة بالرسائل.
  \item البث المباشر الجماعي: «+» في الزاوية اليمنى العلوية > [بث المجموعة]، لبدء بث مباشر يشاهده أعضاء المجموعة فقط، مناسب للتدريب الداخلي والتجمعات العائلية واختبار المنتجات.
  \item طي المجموعة: اسحب اسم المجموعة يسارًا في قائمة المجموعات > [طي]، لجمع المجموعات غير المهمة في «المجموعات المطوية» مع التنبيه فقط عند الإشارة إليك.
  \item الإشارة إلى الجميع من المالك/المشرف فقط: عندما يتجاوز عدد أعضاء المجموعة 500، لا يمكن للأعضاء العاديين الإشارة إلى الجميع، لتجنب الإزعاج.
\end{itemize}

\booksubsection{إدارة دقيقة لصلاحيات Moments:}

\begin{itemize}
  \item إمكانية الرؤية حسب المجموعات: عند النشر انقر على «من يمكنه الرؤية»، واختر «مرئي جزئيًا» أو «إخفاء عن»، اعتمادًا على مجموعات الوسوم المعدة مسبقًا (مثل «الزملاء» و«العائلة» و«العملاء») للنشر الدقيق.
  \item نطاق عرض Moments للأصدقاء: [أنا] > [الإعدادات] > [صلاحيات الأصدقاء] > [Moments] > [نطاق السماح للأصدقاء بعرض Moments]، يمكنك اختيار «آخر ثلاثة أيام/شهر/ستة أشهر/الكل»، ويُوصى بـ«ستة أشهر» لتحقيق التوازن بين الخصوصية والتواصل.
  \item حظر منشور Moments معين: اضغط مطولًا على منشور > [لا أرى Moments الخاص به/ها] أو [إلغاء المتابعة]، ولن يشعر الطرف الآخر.
  \item إدارة إعلانات Moments: [أنا] > [الإعدادات] > [حول WeChat] > [المعلومات الشخصية والصلاحيات] > [إدارة الإعلانات المخصصة]، يمكنك إيقاف التوصيات المخصصة المستندة إلى السلوك.
  \item عشر منشورات للغرباء: [صلاحيات الأصدقاء] > [Moments] > [السماح للغرباء بعرض عشر منشورات]، أوقفها لمنع الغرباء من تصفح منشوراتك السابقة.
\end{itemize}

\booksection{رابعًا: مقارنة مكالمات الفيديو}

تختلف أدوات مكالمات الفيديو في تركيزها، واختيار الأداة المناسبة حسب الموقف يجعل التواصل أكثر سلاسة.

\begin{itemize}
  \item فيديو WeChat: تغطية واسعة، مناسب للجميع. مثالي للتواصل اليومي مع الأهل والأصدقاء والتواصل السريع. يدعم حتى 20 شخصًا (عبر تطبيق «فيديو المجموعة» المصغّر)، وعيوبه أن جودة الصورة تنخفض بوضوح عند ضعف الشبكة.
  \item FaceTime: خاص بمنظومة Apple، جودة ممتازة ومشفر من طرف إلى طرف. منذ iOS 15 يدعم دعوة مستخدمي Android/Web عبر «رابط»، لكن المضيف يجب أن يكون جهاز Apple. مناسب للتواصل العائلي عبر الأجهزة المختلفة.
  \item Zoom: الخيار الأول للمؤتمرات الدولية، يدعم اجتماعات مجانية حتى 100 شخص لمدة 40 دقيقة. ميزات متكاملة مثل مشاركة الشاشة والخلفيات الافتراضية وغرف النقاش، مناسبة للتعاون عبر الحدود/عن بعد.
  \item Tencent Meeting: الخيار الأول للمكاتب في الصين، متكامل بعمق مع WeChat، يدخل بمسح الرمز. يدعم اجتماعات 300 شخص (النسخة المدفوعة)، مع تسجيل سحابي وترجمة فورية ومحاضر اجتماعات قوية، مناسبة للتدريب الداخلي والاجتماعات الكبيرة.
  \item توصيات الاستخدام: للتواصل اليومي فردي اختر WeChat؛ بين مستخدمي Apple اختر FaceTime؛ للاجتماعات الرسمية اختر Zoom أو Tencent Meeting (في الصين Tencent، وعبر الحدود Zoom).
\end{itemize}

\booksubsection{متطلبات النطاق الترددي واستهلاك البيانات:}

\begin{itemize}
  \item مكالمة فيديو WeChat: تحتاج حوالي 500Kbps-1Mbps، وجودة قياسية تستهلك حوالي 3-5 ميجابايت في الدقيقة، والدقة العالية حوالي 8-12 ميجابايت.
  \item FaceTime: الدقة القياسية حوالي 1Mbps، و720p حوالي 2Mbps، و1080p تحتاج 4Mbps فأكثر. الاستهلاك حوالي 5-15 ميجابايت في الدقيقة.
  \item Zoom/Tencent Meeting: 720p حوالي 1.5-3Mbps، و1080p تحتاج 3-5Mbps. الاستهلاك حوالي 300-800 ميجابايت في الساعة، والمضيف في الاجتماعات متعددة الأشخاص يستهلك أكثر.
  \item توصية: استخدم WiFi قدر الإمكان في مكالمات الفيديو الطويلة؛ وتحت شبكة البيانات المتنقلة يمكن خفض الجودة يدويًا (بإيقاف «الفيديو عالي الدقة» في إعدادات التطبيق) لتوفير البيانات.
\end{itemize}

\booksubsection{حيل تحسين جودة مكالمات الفيديو:}

\begin{itemize}
  \item الإضاءة: كن ظهرك نحو النافذة ووجهك نحو مصدر الضوء (الضوء الطبيعي أفضل)، وتجنب الإضاءة الخلفية التي تجعل الوجه داكنًا.
  \item الزاوية: اجعل كاميرا الهاتف أعلى قليلًا من مستوى العينين لتجنب الذقن المزدوج؛ وفي الوضع الأفقي أمسك الهاتف بكلتا اليدين لمزيد من الثبات.
  \item الصوت: استخدم السماعات لتجنب الصدى والصفير؛ ولا تفعّل وضع «مكبر الصوت» في البيئات الهادئة.
  \item الشبكة: أعطِ WiFi الأولوية؛ وفي 4G/5G اقترب من النافذة لتجنب ضعف الإشارة؛ وفي الاجتماعات متعددة الأشخاص يمكنك إيقاف الفيديوهات الأخرى والاحتفاظ بالمتحدث فقط.
  \item التمويه الخلفي: يدعم كلا من Tencent Meeting وZoom التمويه بالذكاء الاصطناعي لإخفاء الخلفية الفوضوية. ارتدِ ملابس ذات لون واحد، وتجنب الألوان القريبة من الخلفية لمنع «انكشف الشاشة الخضراء».
\end{itemize}

\begin{tcbinfo}{نصيحة}
عند بطء الشبكة، أوقف الفيديو أولًا وأبقِ الصوت، فهو أكثر أدبًا من قطع المكالمة مباشرة. قبل الاجتماع المهم يمكنك استخدام ميزة «فحص الشبكة» المدمجة في Tencent Meeting لاختبار النطاق الترددي مسبقًا.

\end{tcbinfo}

\booksection{خامسًا: البريد الإلكتروني: الإعداد والاستخدام على الهاتف}

\begin{itemize}
  \item إضافة حساب بريد: في iOS من الإعدادات > البريد > الحسابات > إضافة حساب؛ وفي Android يوجد عادةً تطبيق «البريد الإلكتروني».
  \item نصائح الاستخدام: اضبط طرق إشعار مختلفة لكل حساب بريد؛ واستخدم إيماءات السحب للأرشفة أو الحذف أو الوسم السريع.
\end{itemize}

\booksubsection{الفرق بين بروتوكولات IMAP وPOP3 وExchange:}

\begin{itemize}
  \item POP3: بعد تنزيل الرسائل من الخادم يُحذف الأصل افتراضيًا ويُحفظ محليًا فقط. مناسب لاستخدام جهاز واحد والقراءة دون اتصال. عيوبه عدم ظهور الرسائل السابقة عند تغيير الجهاز، ولم يعد موصى به.
  \item IMAP: تُحفظ الرسائل دائمًا على الخادم، والنسخة المحلية مجرد ذاكرة مؤقتة. مزامنة متسقة عبر الأجهزة (هاتف/كمبيوتر/جهاز لوحي) — إذا قُرئت في مكان واحد تُقرأ في الكل. وهو البروتوكول الموصى به حاليًا.
  \item Exchange (بما في ذلك ActiveSync): بروتوكول Microsoft، يزامن بالإضافة إلى البريد التقويم وجهات الاتصال والمهام. الخيار الأول لبريد الشركات. يدعمه Outlook وNetEase Enterprise Mail وTencent Enterprise Mail.
  \item توصية الاختيار: للبريد الشخصي استخدم IMAP؛ ولبريد الشركة فضل Exchange ثم IMAP؛ وتجنب POP3 قدر الإمكان إلا إذا كانت هناك حاجة خاصة للأرشفة.
\end{itemize}

\booksubsection{طريقة إعداد بريد الشركات:}

\begin{enumerate}
  \item iOS: [الإعدادات] > [البريد] > [الحسابات] > [إضافة حساب] > اختر «Microsoft Exchange» أو «أخرى».
  \item Android: في تطبيق «البريد الإلكتروني» المدمج أو تطبيق Outlook اختر «إضافة حساب»، وأدخل عنوان البريد الكامل وكلمة المرور.
  \item إعداد الخادم (بمثال Exchange): عنوان الخادم عادةً «mail.نطاق-الشركة.com» أو «outlook.office365.com»، والاسم النطاق هو اسم الشركة المختصر.
  \item إذا ظهرت رسالة «تعذر التحقق من الخادم»، أدخل يدويًا: خادم الوارد (IMAP) المنفذ 993 مع تشفير SSL؛ وخرج الخادم (SMTP) المنفذ 465 أو 587 مع تشفير TLS.
  \item بعد تسجيل الدخول يُنصح بتفعيل «الدفع» (Push) في [إعدادات الحساب] لاستقبال الرسائل فورًا، مع إيقاف «معاينة البريد» لحماية الأسرار التجارية.
\end{enumerate}

\begin{tcbinfo}{ملاحظة}
عادةً يحتاج بريد الشركات إلى تفعيل «كلمة مرور تفويض العميل» أو «رمز تفويض IMAP/SMTP» من المسؤول، ولا يمكن استخدام كلمة مرور تسجيل الدخول على الويب مباشرة. عند فشل الدخول أنشئ رمز التفويض من إعدادات بريد الويب أولًا ثم أدخله في الهاتف.

\end{tcbinfo}

\begin{tcbpractice}{تطبيق عملي: نسخ سجلات محادثات WeChat إلى الكمبيوتر}
\begin{enumerate}
  \item ابدأ النسخ الاحتياطي على الكمبيوتر: افتح WeChat على الكمبيوتر > القائمة في الزاوية اليسرى السفلية > [النقل والنسخ الاحتياطي] > [النسخ الاحتياطي والاستعادة].
  \item اختر النسخ الاحتياطي إلى الكمبيوتر: انقر على «النسخ الاحتياطي لسجلات المحادثة إلى الكمبيوتر».
  \item حدد على الهاتف محتوى النسخ الاحتياطي، وانتظر اكتمال العملية.
\end{enumerate}

\end{tcbpractice}
